% ============================================================
% Kartu faktor — di-\input dari Subbab 4.11.1 (Bab4.tex).
% Gaya kartu (kartufaktor, kotakekspresi, fasebadge, barisfield)
% dan paket (booktabs, tcolorbox) didefinisikan di preamble skripsi.tex.
% Mengikuti format Lampiran C QuantaAlpha (Factor Identity + Backtest).
% ============================================================

% Kartu-kartu ini TIDAK dibungkus \begin{figure}: kartufaktor bersifat
% breakable (bisa melebihi satu halaman), sedangkan sebuah float wajib muat
% dalam satu halaman -- kombinasi keduanya membuat pdflatex diam-diam
% memangkas isi yang melebihi tinggi halaman (lihat "Float too large").
% Nomor Gambar tetap didapat lewat \captionof{figure} (paket caption).
\begin{kartufaktor}{small\_cap\_stocks\_with\_unusually\_1}

\barisfield{Factor ID:}{\texttt{0ac0ab9b1b9323f8}}
\barisfield{Trajectory ID:}{\texttt{6e2654008cae}}
\barisfield{Ronde Evolusi:}{Ronde 0}
\barisfield{Fase Evolusi:}{\fasebadge{Original}}
\barisfield{Mode Komunikasi:}{\texttt{text}}
\barisfield{Arah Eksplorasi:}{\textit{short-term reversal after abnormally
  high-volume days in small-cap stocks}}

\vspace{2mm}
\textbf{Hipotesis:}\par
\textit{``small-cap stocks with unusually high volume on a day are followed by
negative cross-sectional returns, reflecting liquidity-induced mean
reversion.''}

\vspace{2.5mm}
\textbf{Ekspresi Faktor:}
\begin{kotakekspresi}
RANK(\$volume) * (TS\_RANK(\$return, 1) ? -1 : 1)
\end{kotakekspresi}

\textbf{Formulasi Matematis:}
\begin{equation*}
  f_t \;=\; \operatorname{RANK}(V_t)\times
  \begin{cases}
    -1, & \text{jika } \operatorname{TS\_RANK}(r_t,\,1)\neq 0,\\[1mm]
    \phantom{-}1, & \text{lainnya,}
  \end{cases}
\end{equation*}
dengan $\operatorname{RANK}(\cdot)$ peringkat persentil \textit{cross-sectional},
$V_t$ volume harian, $r_t$ \textit{return} harian, dan
$\operatorname{TS\_RANK}(\cdot,1)$ peringkat deret-waktu berjendela satu hari.

\vspace{2mm}
\textbf{Diagnosis Agen \textit{Feedback}:}\par
\textit{``strong, valid mechanism but implementation is shallow''} ---
kelima ekspresi saudara didiagnosis
\textit{noise} (\textit{missing normalization}, \textit{wrong window},
\textit{shallow composition}). Putusan hipotesis:
\textit{``SUPPORTS --- the mechanism is valid but the implementation is weak.''}

\end{kartufaktor}
\captionof{figure}{Kartu faktor \textit{trajectory} terbaik batch komparasi
         (mode \texttt{text}, fase \textit{original}), mengikuti format
         Lampiran~C QuantaAlpha.}
\label{fig:kartu-faktor-terbaik}

\vspace{4mm}
\begin{kartufaktor}{Metrik Backtest --- Trajectory \texttt{6e2654008cae}}

\begin{center}
\begin{tabular}{lc}
\toprule
\textbf{Metrik} & \textbf{Nilai} \\ \midrule
RankIC \textit{standalone} (\texttt{FactorIC\_mean}) & $\mathbf{0{,}0449}$ \\
ICIR \textit{standalone}                             & $0{,}4310$ \\ \midrule
IC (model gabungan)        & $0{,}0165$ \\
RankIC (model gabungan)    & $0{,}0359$ \\ \midrule
ARR (\textit{excess, w/ cost})  & $9{,}60\%$ \\
IR (\textit{w/ cost})           & $0{,}805$ \\
MDD (\textit{excess, w/ cost})  & $-8{,}35\%$ \\ \bottomrule
\end{tabular}
\end{center}

\vspace{2mm}
\textbf{Statistik Rinci:}
\begin{center}
\small
\begin{tabular}{lc@{\hspace{8mm}}lc}
\toprule
\textbf{Metrik} & \textbf{Nilai} & \textbf{Metrik} & \textbf{Nilai} \\ \midrule
Daily Excess Return (w/o cost) & $0{,}0501\%$ & Daily Excess Return (w/ cost) & $0{,}0403\%$ \\
Excess Return Std              & $0{,}77\%$   & Turnover (FFR)                & $100\%$ \\
L2 Train Loss                  & $0{,}9938$   & L2 Valid Loss                 & $0{,}9969$ \\
\bottomrule
\end{tabular}
\end{center}

\end{kartufaktor}
\captionof{figure}{Metrik \textit{backtest out-of-sample} (2022--2025) \textit{trajectory}
         terbaik; RankIC \textit{standalone} bebas model, metrik lain berasal
         dari \textit{backtest} model gabungan LightGBM.}
\label{fig:metrik-backtest-terbaik}

\vspace{4mm}
\begin{kotakrespon}{Respon Mentah Agen --- Trajectory \texttt{6e2654008cae} (mode \texttt{text})}

\textbf{Arah Eksplorasi:} \textit{short-term reversal after abnormally
high-volume days in small-cap stocks}

\vspace{2mm}
\textbf{Agen Proposal:}\par
\textit{``observation: small-cap stocks with abnormally high volume on a
given day tend to show reversal in the next day's cross-sectional
returns.}

\textit{driver: this is a liquidity-driven reversal, where high volume
indicates excessive buying or selling pressure, leading to mean
reversion. the signal is derived from volume and return columns, falls
within a short horizon band, and predicts negative returns across the
cross-section.}

\textit{HYPOTHESIS: small-cap stocks with unusually high volume on a day
are followed by negative cross-sectional returns, reflecting
liquidity-induced mean reversion.''}

\vspace{2mm}
\textbf{Agen Design} (\textit{shortlisted functions}):

\noindent\textbf{1.} \texttt{TS\_ZSCORE(\$volume, 1)} --- computes the
z-score of volume over the past 1 day, measuring how unusual the volume
is compared to other stocks, indicating liquidity pressure.\par
\noindent\textbf{2.} \texttt{TS\_MEAN(\$return, 1)} --- computes the mean
of returns over the past 1 day, the cross-sectional return average
targeted by the reversal.\par
\noindent\textbf{3.} \texttt{TS\_RANK(\$return, 1)} --- ranks returns
across the cross-section today, helping identify stocks with negative
returns.\par
\noindent\textbf{4.} \texttt{RANK(\$volume)} --- ranks volume across the
cross-section today, identifying stocks with unusually high volume.\par
\noindent\textbf{5.} \texttt{TS\_ZSCORE(\$return, 1) ? -1 : 1} --- a
binary score based on the z-score of returns; negative values indicate
negative returns.\par
\noindent\textbf{6.} \texttt{\$return ? -1 : 1} --- a simple binary score
based on the return itself, indicating negative or positive returns.

\vspace{2mm}
\textbf{Agen Construct} (enam ekspresi faktor kandidat):

\noindent\textbf{1.} \texttt{VolumeZScoreReturn}: \texttt{TS\_ZSCORE(\$volume,
1) * (TS\_MEAN(\$return, 1) ? -1 : 1)} --- multiplies the z-score of
volume (indicating liquidity pressure) with a binary score of
cross-sectional returns (negative values indicate negative
returns).\par
\noindent\textbf{2.} \texttt{VolumeRankReturn}: \texttt{RANK(\$volume) *
(TS\_RANK(\$return, 1) ? -1 : 1)} --- multiplies the rank of volume
(higher ranks indicate higher volume) with a binary score of returns
(negative values indicate negative returns). \textit{[ekspresi terpilih;
bandingkan Gambar~\ref{fig:kartu-faktor-terbaik}]}\par
\noindent\textbf{3.} \texttt{VolumeZScoreReturnGate}:
\texttt{(TS\_ZSCORE(\$volume, 1) $>$ 1) ? TS\_MEAN(\$return, 1) : 0} ---
gates the cross-sectional return mean based on the z-score of volume
(only active for stocks with unusually high volume).\par
\noindent\textbf{4.} \texttt{VolumeZScoreReturnQuantile}:
\texttt{TS\_QUANTILE(\$return, 0.1, 1) * TS\_ZSCORE(\$volume, 1)} ---
combines the 10th percentile of returns (indicating negative returns)
with the z-score of volume to highlight stocks with both high volume and
negative returns.\par
\noindent\textbf{5.} \texttt{VolumeReturnCrossRank}: \texttt{TS\_RANK(\$return,
1) * (TS\_ZSCORE(\$volume, 1) ? -1 : 1)} --- combines the rank of
returns (identifying negative returns) with a binary score of volume
z-scores to emphasize stocks with high volume and negative returns.\par
\noindent\textbf{6.} \texttt{VolumeReturnZScoreGate}:
\texttt{(TS\_ZSCORE(\$volume, 1) $>$ 1) ? TS\_ZSCORE(\$return, 1) : 0} ---
gates the z-score of returns based on the z-score of volume (only active
for stocks with unusually high volume).

\textbf{Agen Feedback:}\par
\textbf{Observations:} \textit{The SOTA is determined by [A] standalone
RankIC+ICIR, where FactorIC\_mean has RankIC=0.0449 and ICIR=0.4310. [B]
combined model only provides context and does not decide SOTA.}

\textbf{Factor Diagnosis:} \textit{small\_cap\_stocks\_with\_unusually\_1:
strong, valid mechanism but implementation is shallow.
small\_cap\_stocks\_with\_unusually\_2: noise, missing normalization.
small\_cap\_stocks\_with\_unusually\_3: noise, wrong window.
small\_cap\_stocks\_with\_unusually\_4: noise, shallow composition.
small\_cap\_stocks\_with\_unusually\_5: noise, missing normalization.}

\textbf{Feedback for Hypothesis:} \textit{SUPPORTS --- the mechanism is
valid but the implementation is weak (refine the expression). The
framework still has room for refinement.}

\textbf{Replace Best Result:} \textit{yes}

\end{kotakrespon}
\captionof{figure}{Respon mentah keempat agen \textit{front-end} (proposal,
         design, construct, feedback) untuk \textit{trajectory} yang sama,
         mengikuti format kotak studi kasus Lampiran D LatentMAS
         \citep{latentmas}; teks dipertahankan dalam bahasa asli keluaran
         model.}
\label{fig:respon-mentah}
